\documentclass[11pt]{article}

\usepackage[spanish]{babel}
\usepackage[utf8]{inputenc}
\usepackage[T1]{fontenc}
\usepackage{amsmath,amssymb}
\usepackage{booktabs}
\usepackage{geometry}
\usepackage{enumitem}
\usepackage{xcolor}
\usepackage{tcolorbox}
\usepackage{hyperref}

\geometry{margin=1in}

\hypersetup{
    colorlinks=true,
    linkcolor=blue
}

\title{\textbf{Ejercicio Integrador de Stress Testing}\\
\vspace{0.3cm}
\large De los Factores de Riesgo al Impacto sobre el Portafolio Previsional}
\author{}
\date{}

\begin{document}

\maketitle


\section*{Objetivo}

Comprender cómo un shock macroeconómico se transforma en una distribución de posibles retornos y pérdidas para un fondo de pensiones.

Al finalizar el ejercicio, el estudiante será capaz de:

\begin{itemize}
\item Identificar factores de riesgo relevantes.
\item Construir escenarios macroeconómicos consistentes.
\item Estimar la sensibilidad del portafolio a dichos factores.
\item Simular trayectorias futuras del portafolio.
\item Cuantificar pérdidas potenciales mediante VaR y Expected Shortfall.
\end{itemize}

\begin{tcolorbox}[colback=blue!5,colframe=blue!50!black,title=Pregunta central]
¿Cómo se transforma un shock macroeconómico en una distribución de pérdidas para un fondo de pensiones?
\end{tcolorbox}

\section{El problema que queremos resolver}

Imagine que usted forma parte del Comité de Riesgos de una AFP.

Durante una reunión, un miembro del comité pregunta:

\begin{quote}
``¿Qué ocurriría con la rentabilidad del fondo si el crecimiento se contrae, la inflación aumenta, las tasas de interés aumentaran, el tipo de cambio se depreciara y el riesgo soberano aumentara simultáneamente?''
\end{quote}

Nadie puede predecir el futuro con certeza.

Sin embargo, sí es posible construir escenarios plausibles y cuantificar sus consecuencias potenciales.

Eso es precisamente un \textbf{stress test}.

\section{Paso 0. Construcción del mapa mental del problema}

Antes de estimar cualquier modelo debemos comprender los mecanismos económicos que conectan los factores macroeconómicos con el desempeño del portafolio.

\subsection*{Pregunta}

¿Cómo afecta cada factor al portafolio?

Complete la siguiente tabla:

\begin{center}
\begin{tabular}{lll}
\toprule
Factor & Canal de transmisión & Efecto esperado \\
\midrule
Tasa de interés & Precio de bonos & Negativo \\
Tipo de cambio & Activos externos & Positivo \\
Mercado accionario & Valor de acciones & Positivo \\
Spread soberano & Riesgo país & Negativo \\
\bottomrule
\end{tabular}
\end{center}

\subsection*{Idea central}

\begin{center}
\begin{verbatim}
Shock macroeconómico
          ↓
Mercados financieros
          ↓
Portafolio
          ↓
Ganancias o pérdidas
\end{verbatim}
\end{center}

El stress testing comienza entendiendo esta cadena causal.

\section{Paso 1. ¿Cómo se mueven los factores de riesgo?}

Suponga que la tasa de interés aumenta hoy.

\begin{itemize}
\item ¿Qué ocurrirá con el tipo de cambio?
\item ¿Qué ocurrirá con el mercado accionario?
\item ¿Qué ocurrirá con el spread soberano?
\end{itemize}

Los factores de riesgo interactúan entre sí.

Para modelar dichas interacciones estimaremos un modelo VAR.

\subsection*{Tarea}

Estime un VAR utilizando:

\begin{itemize}
\item Tasa de interés.
\item Tipo de cambio.
\item Índice accionario.
\item Spread soberano.
\end{itemize}

\subsection*{Producto esperado}

Presentar:

\begin{itemize}
\item Selección de rezagos.
\item Funciones impulso-respuesta.
\item Interpretación económica de los resultados.
\end{itemize}

\subsection*{Interpretación}

El VAR es una máquina de propagación de shocks.

\begin{center}
\begin{verbatim}
Shock inicial
      ↓
VAR
      ↓
Trayectoria futura de los factores
\end{verbatim}
\end{center}

\section{Paso 2. Construcción de escenarios}

Ahora utilizaremos el VAR para responder:

\begin{quote}
``Si ocurre un evento adverso, ¿cómo evolucionarán los factores de riesgo?''
\end{quote}

\subsection*{Escenario Base}

La economía continúa su trayectoria esperada.

\subsection*{Escenario Adverso}

Ejemplo:

\begin{itemize}
\item Crecimiento cae a -2\%
\item Inflación aumenta 300 pb.
\item Aumento de 300 pb en tasas.
\item Depreciación de 10\%.
\item Aumento de 200 pb del riesgo país.
\end{itemize}

\subsection*{Escenario Severo}

Ejemplo:

\begin{itemize}
\item Crecimiento cae a -6%
\item Inflación aumenta 550 pb.
\item Aumento de 500 pb en tasas.
\item Depreciación de 20\%.
\item Aumento de 400 pb del riesgo país.
\end{itemize}

\subsection*{Producto esperado}

Graficar la evolución proyectada de:

\begin{itemize}
\item Tasas de interés.
\item Tipo de cambio.
\item Índice accionario.
\item Spread soberano.
\end{itemize}

\subsection*{Pregunta de reflexión}

¿Cuál escenario parece más perjudicial para el portafolio?

¿Por qué?

\section{Paso 3. ¿Cómo responde el portafolio?}

Hasta ahora hemos modelado los factores.

Ahora debemos conectar dichos factores con la rentabilidad del portafolio.

\subsection*{Modelo de factores}

\[
R_t =
\alpha
+
\beta_1 Tasa_t
+
\beta_2 FX_t
+
\beta_3 Equity_t
+
\beta_4 Spread_t
+
\varepsilon_t
\]

donde:

\begin{itemize}
\item $R_t$ representa la rentabilidad del portafolio.
\item Los $\beta$ representan sensibilidades.
\end{itemize}

\subsection*{Interpretación intuitiva}

Si:

\[
\beta_{Equity}=0.40
\]

entonces un aumento de un punto porcentual en el mercado accionario incrementa la rentabilidad esperada del portafolio en aproximadamente 0.4 puntos porcentuales.

\subsection*{Producto esperado}

Interpretar:

\begin{itemize}
\item Signo de los coeficientes.
\item Magnitud económica.
\item Significancia estadística.
\end{itemize}

\section{Paso 4. Construcción de la trayectoria esperada del portafolio}

Ya disponemos de:

\begin{enumerate}
\item Escenarios para los factores.
\item Sensibilidades del portafolio.
\end{enumerate}

Ahora combinamos ambas piezas.

\subsection*{Idea}

\begin{center}
\begin{verbatim}
Escenario de factores
            +
Modelo de factores
            ↓
Retorno esperado
\end{verbatim}
\end{center}

\subsection*{Resultado}

Obtener:

\[
\widehat{R}_{t+h}
\]

para:

\begin{itemize}
\item Baseline.
\item Adverso.
\item Severo.
\end{itemize}

\subsection*{Pregunta}

¿Es suficiente conocer el retorno esperado?

La respuesta es no.

La realidad contiene incertidumbre.

\section{Paso 5. Incorporando incertidumbre}

Imagine que dos economías enfrentan exactamente el mismo escenario.

Aun así, los resultados pueden ser distintos.

¿Por qué?

Porque existen shocks impredecibles.

\subsection*{Idea}

\[
R_t
=
\widehat{R}_t
+
\varepsilon_t
\]

donde:

\begin{itemize}
\item $\widehat{R}_t$ es el retorno esperado.
\item $\varepsilon_t$ representa incertidumbre residual.
\end{itemize}

\subsection*{Implementación}

\begin{enumerate}
\item Estimar un modelo GARCH sobre los residuos históricos.
\item Simular residuos futuros.
\item Agregarlos a la trayectoria esperada.
\end{enumerate}

\subsection*{Flujo}

\begin{center}
\begin{verbatim}
Retorno esperado
         +
Residuos simulados
         ↓
Miles de trayectorias posibles
\end{verbatim}
\end{center}

\subsection*{Tarea}

Generar al menos 10,000 simulaciones Monte Carlo.

\section{Paso 6. Del escenario a la pérdida}

Ahora disponemos de una distribución completa de resultados posibles.

Necesitamos resumirla mediante métricas de riesgo.

\subsection*{Value-at-Risk (VaR)}

Pregunta:

\begin{quote}
``¿Cuál es el retorno que sólo sería superado en el peor 1\% de los casos?''
\end{quote}

\[
VaR_{99}
\]

\subsection*{Expected Shortfall (ES)}

Pregunta:

\begin{quote}
``Si entramos en el peor 1\% de los casos, ¿cuál sería el retorno promedio?''
\end{quote}

\[
ES_{99}
\]

\subsection*{Producto esperado}

\begin{center}
\begin{tabular}{lrrrr}
\toprule
Escenario & Retorno Esperado & VaR 95\% & VaR 99\% & ES 99\% \\
\midrule
Baseline & & & & \\
Adverso & & & & \\
Severo & & & & \\
\bottomrule
\end{tabular}
\end{center}

\section{Presentación al Comité de Inversiones}

Suponga que dispone de únicamente dos minutos para presentar sus resultados.

Responda:

\begin{enumerate}
\item ¿Cuál es el principal factor de riesgo?
\item ¿Qué escenario genera la mayor pérdida?
\item ¿Cuál es la magnitud de la pérdida extrema?
\item ¿Qué medidas de mitigación recomendaría?
\end{enumerate}

\section{La idea fundamental del ejercicio}

\begin{center}
\begin{verbatim}
Narrativa económica
          ↓
Factores de riesgo
          ↓
VAR
          ↓
Escenarios
          ↓
Modelo de factores
          ↓
Retorno esperado
          ↓
GARCH + Monte Carlo
          ↓
Distribución de retornos
          ↓
VaR y Expected Shortfall
          ↓
Decisión de gestión de riesgo
\end{verbatim}
\end{center}

\begin{tcolorbox}[colback=green!5,colframe=green!50!black,title=Mensaje Final]
Un stress test no consiste en producir números.

Consiste en transformar una narrativa económica en una medida cuantitativa del riesgo que permita apoyar la toma de decisiones.
\end{tcolorbox}

\end{document}